\scanpage{68}
\mspage{67}
% E: text|upper left|au départ pour repérer une op.
au départ, pour repérer une op. de générisation $D'\longrightarrow D$ par un ``arc de \deleted{\unsure{spéciali}} générisation'' sur $\partial\Delta_{D'}$; on prend maintenant, non la rive qui ``fait face'' à celle de $b$, mais par contre celle qui en résulte par transport parallèle. \deleted{\unsure{Je pensais}} \inserted{Quant à} l'isom. de transition $\underline\omega_{\Delta_D}\mathbin{-}\underline\omega_{\Delta_{D'}}$ pour les orientations, faisant comme \unsure{d'hab.}
% E: diagram|left middle|omega orientations and transport
\diagram{\begin{tikzcd}[column sep=large,row sep=large]
\omega\ \text{or. de }\Delta_D \arrow[r,leftrightarrow] \arrow[dr,leftrightarrow,bend right=45,"\text{transport}"'] & \omega_D\ \text{or. de }D & \\
& \omega'\ \text{or. de }\Delta_{D'} \arrow[r,leftrightarrow] & \omega_{D'}\ \text{or. de }D',
\end{tikzcd}}
% E: text|left below diagram|on trouve à nouveau
on trouve à nouveau que $\omega_{D'}$ est \uline{opposée} à l'orientation \deleted{\unsure{de $D$}} déduite de $\omega_D$ par transport parallèle!

Pendant ces deux dernières semaines, c'est avec cette opération de recollage que j'ai travaillé. Avant, j'avais pris \deleted{\unsure{même}} l'\uline{opposé} iso. de transition pour les orientations, \unsure{le|Ce} transport parallèle --- \deleted{\unsure{de la face}} et comme \deleted{\unsure{même}} recollage, si je me rappelle bien, le recollage
% E: text|upper right|1 degree de b avec a'
$1^\circ$) de $b$ avec $a'$ --- je croyais avoir interprété correctement la construction de recollement de $\unsure{4|1}^\circ$), mais je m'étais mélangé les pinceaux de \uline{deux} façons, et \unsure{ensemble}:

1) Je m'étais trompé sur les isom. de transition pour les orientations, et

2) je m'étais trompé également sur l'identification correcte d'une op. de générisation \inserted{pour $D'$} par un arc sur $\partial\Delta_{D'}$ --- en prenant \unsure{en symbole la bonne|en somme la même} celle qui \unsure{en est} \deleted{\unsure{est}} déduit de $\unsure{b}$ par transport parallèle, tandis que c'est l'autre \deleted{qu'il} (l'antipodique, qu'il fallait prendre).

Pour résumer, les quatre possibilités sont indiquées dans la liste suivante
